% P05 — Accountable but Unlinkable Agent Identity — arXiv preprint
% Build: tectonic -Z shell-escape main.tex  (XeTeX engine; fetches packages)
% STATUS: working draft. The Results section reads generated LaTeX from
% ../results/; do not type a number into this file by hand.
%
% Preamble, colour palette, \aaimark and the cover-page TikZ block are copied
% from ../../ov-poc-standard/paper/main.tex so the papers in this series look
% like one series. One deliberate difference, already established by the
% sibling P01 paper in ../../parallax/paper/main.tex: that paper loads
% `minted`, which needs a pygmentize binary on the build host. This paper has
% no syntax-highlighted listings and the build host here has no pygmentize
% installed, so minted is not loaded. `listings` is configured identically in
% case that changes.
\documentclass[11pt]{article}

\usepackage[margin=1in]{geometry}
\usepackage{tikz}
\usepackage{graphicx}
\usepackage{booktabs}
\usepackage{array}
\usepackage{amsmath, amssymb}
\usepackage{listings}
\definecolor{codepanel}{HTML}{F7F5F1}
\usepackage{xcolor}
\usepackage{colortbl}
\usepackage[hidelinks]{hyperref}
\usepackage{enumitem}
\usepackage{microtype}
\usepackage{caption}
\usepackage{siunitx}
\usepackage{amsthm}
\newtheorem{definition}{Definition}
\newtheorem{theorem}{Theorem}
\newtheorem{proposition}{Proposition}
\usepackage[section]{placeins}

% Advanced AI Society brand palette
\definecolor{aailime}{HTML}{CFFF04}
\definecolor{aailimetint}{HTML}{F2FFB8}
\definecolor{aaiink}{HTML}{0A0A0A}
\definecolor{aaired}{HTML}{FFE3E4}
\definecolor{aaineutral}{HTML}{E7E3DD}
\definecolor{aaipaper}{HTML}{F0EDEA}
\definecolor{aaimuted}{HTML}{8F8A82}
\definecolor{aaidarklime}{HTML}{7A9900}

\lstset{basicstyle=\ttfamily\footnotesize, breaklines=true, frame=single,
        framerule=0.3pt, xleftmargin=1em, columns=fullflexible}
\captionsetup{font=small, labelfont=bf}

\newcommand{\occultation}{\texttt{occultation}}
% An underscore inside a code identifier, with a break opportunity after it.
\newcommand{\ub}{\_\allowbreak}

% The Society's bracket mark, drawn natively (lime on the dark band)
\newcommand{\aaimark}[1]{%
\begin{tikzpicture}[x=#1, y=#1, line cap=butt]
  \fill[aailime] (0,0) rectangle (0.16,1);
  \fill[aailime] (0.16,0) rectangle (0.42,0.24);
  \fill[aailime] (0.16,0.76) rectangle (0.42,1);
  \fill[aailime] (1.04,0) rectangle (1.20,1);
  \fill[aailime] (0.78,0) rectangle (1.04,0.24);
  \fill[aailime] (0.78,0.76) rectangle (1.04,1);
\end{tikzpicture}}

\begin{document}

% Generated-macro fact sheets, loaded before first use: the abstract itself
% cites pool and composed-path figures, so these cannot wait for
% \section{Results}. See scripts/regen-results.sh; nothing below this line
% and above \tableofcontents is typed by hand.
% Generated by scripts/regen-results.sh from the desk_study rows of
% results/measured/composed.json (the same field Table 3 reads). These are
% the desk study's published constants, quoted in prose before Table 3
% appears; formatted identically to it so the two cannot silently diverge.
\newcommand{\deskbudget}{\SI{15.00}{\milli\second}}
\newcommand{\deskpresentnaive}{\SI{13.80}{\milli\second}}
\newcommand{\deskverifynaive}{\SI{4.20}{\milli\second}}
\newcommand{\desktotalnaive}{\SI{18.00}{\milli\second}}
\newcommand{\deskovernaive}{\SI{3.00}{\milli\second}}
\newcommand{\desktotalcached}{\SI{14.90}{\milli\second}}
\newcommand{\deskheadroomcached}{\SI{100}{\micro\second}}
\newcommand{\deskverifycached}{\SI{1.10}{\milli\second}}
\newcommand{\deskpresentpooled}{\SI{2.80}{\milli\second}}
\newcommand{\desktotalpooled}{\SI{3.90}{\milli\second}}
\newcommand{\deskheadroompooled}{\SI{11.10}{\milli\second}}
\newcommand{\deskcachedsavingpct}{\SI{74}{\percent}}
\newcommand{\deskpooledreductionpct}{\SI{80}{\percent}}

% Generated by scripts/regen-results.sh from
% results/deterministic/pool-burst-5000.json, the same burst
% tests/acceptance.rs's acceptance test 2 exercises. A seeded simulation,
% hence deterministic; do not hand-edit.
\newcommand{\poolrequests}{\num{10000}}
\newcommand{\poolserved}{\num{89}}
\newcommand{\poolstalled}{\num{9911}}
\newcommand{\poolamplification}{\num{82591}}
\newcommand{\poolwindow}{\SI{2.0}{\second}}
\newcommand{\pooloverloadfactor}{\num{3.50}}
\newcommand{\poolgrowthperrequest}{\num{8.33}}
\newcommand{\poolonsetrequest}{\num{89}}
\newcommand{\pooldistinctblinding}{\num{10000}}
\newcommand{\poolrefillrate}{\num{1429}}
\newcommand{\poolburstrate}{\num{5000}}
\newcommand{\poolrefillratesi}{\SI{1429}{\per\second}}
\newcommand{\poolburstratesi}{\SI{5000}{\per\second}}
\newcommand{\poolrefillcostsi}{\SI{700}{\micro\second}}
\newcommand{\poolonlinecostsi}{\SI{60}{\micro\second}}
\newcommand{\poolreuseevents}{\num{0}}

% Generated by scripts/regen-results.sh from
% results/deterministic/anonymity-drifted.json (the same file Table 5
% reads). Summary figures quoted in prose outside the table.
\newcommand{\anonfleetsize}{\num{120}}
\newcommand{\anonconfigs}{\num{5}}
\newcommand{\anonsingletons}{\num{1}}
\newcommand{\anonsizes}{71, 34, 11, 3, 1}
\newcommand{\anoneffectiveset}{\num{52.73}}
\newcommand{\anonentropy}{\num{1.470} bits}

% Generated by scripts/regen-results.sh from results/measured/composed.json.
% Wall-clock timing: differs per host, not diffed by CI. See
% results/measured/HEADER.txt for the host and toolchain. Do not hand-edit.
\newcommand{\measureddivergencefactor}{\num{14}}
\newcommand{\measureddivergencedirection}{cheaper}
\newcommand{\measuredcachedsavingpct}{\SI{11}{\percent}}
\newcommand{\measuredratiotobaseline}{\num{52}}
\newcommand{\pooledreductionpct}{\SI{99}{\percent}}


% ---------------------------------------------------------------- cover page
\thispagestyle{empty}

\begin{tikzpicture}[remember picture, overlay]
  \fill[aaiink] (current page.north west) rectangle
    ([yshift=-6.35cm]current page.north east);
  \fill[aailime] ([yshift=-6.35cm]current page.north west) rectangle
    ([yshift=-6.53cm]current page.north east);
  \node[anchor=north west, inner sep=0pt]
    at ([xshift=1in, yshift=-0.95cm]current page.north west) {\aaimark{0.62cm}};
  \node[anchor=north west, inner sep=0pt, text width=14.2cm]
    at ([xshift=1in+1.15cm, yshift=-0.98cm]current page.north west)
    {\sffamily\bfseries\footnotesize\color{aailime}%
     ADVANCED AI SOCIETY\; \textcolor{aaimuted}{$\cdot$}\; PROOF-OF-CONTROL
     INITIATIVE\\[3pt]
     {\normalfont\sffamily\footnotesize\color{aaimuted}
      Working draft v0.1 $\cdot$ August 2026 $\cdot$ for co-author review
      prior to submission}};
  \node[anchor=north west, inner sep=0pt, text width=\textwidth]
    at ([xshift=1in, yshift=-2.35cm]current page.north west)
    {\sffamily\bfseries\fontsize{19}{23}\selectfont\color{aaipaper}%
     What Unlinkability Costs at Agent Action Rates: A Composed-Path
     Measurement and the Soundness of Pre-Computation};
  \node[anchor=north west, inner sep=0pt, text width=\textwidth]
    at ([xshift=1in, yshift=-5.55cm]current page.north west)
    {\sffamily\small\color{aaipaper}%
     \textcolor{aailime}{\rule[0.5ex]{1.6em}{2pt}}\hspace{0.6em}%
     The obvious construction misses the budget, and the fix is a security
     question nobody has asked.};
\end{tikzpicture}

\vspace*{5.9cm}

% Compact author block
{\small\noindent
\textbf{Jim Schwoebel}\textsuperscript{1,2}

\smallskip\noindent
{\footnotesize
\textsuperscript{1}\,Quome \quad
\textsuperscript{2}\,Advanced AI Society, \url{https://advancedaisociety.org}%
\footnote{Prepared within the Proof-of-Control Initiative of the Advanced AI
Society. Authorship is limited to contributors to this manuscript; additional
co-authors will be added only upon explicit consent and substantive
contribution.}}}

\vspace{0.9em}
\noindent\textcolor{aailime}{\rule{\textwidth}{1.6pt}}
\vspace{0.2em}

\begin{center}
{\sffamily\bfseries\small ABSTRACT}
\end{center}
\vspace{-0.5em}
\noindent{\small
An agent presents an anonymous credential and a relying party verifies it,
and the standard gives the two of them together one \deskbudget{} action
budget. A desk study checked them one at a time. BBS+ presentation, about
\deskpresentnaive{}: under. Verification, about \deskverifynaive{}: under.
Two halves under budget, so the design fits --- except that the action pays
for both, and the sum is \desktotalnaive{}, over by \deskovernaive{}. With
cached pairing pre-computation the sum is \desktotalcached{}, under the limit
by \deskheadroomcached{}, which this paper's tool reports as
\texttt{MARGINAL} rather than a pass, because \deskheadroomcached{} is not
headroom for the application, the network, or writing the evidence record.
We built \occultation{}, a Rust tool that assesses presentation and
verification as one composed path rather than two separate ones, over real
BBS+ signatures on BLS12-381 and a real Ed25519 baseline. Our hardware is
much faster --- the composed path costs \measureddivergencefactor{}$\times$
less here than the desk study's published figures imply, or
\measureddivergencefactorexclescrow{}$\times$ once the modelled escrow tag is
excluded from both sides --- and none of that touches the mistake, because
the mistake is in the addition and not in the silicon. What the magnitude
would hide is that the \emph{shape} of the correction survives the hardware
change while the shape of an unrelated optimization does not: cached pairings
save \deskcachedsavingpct{} of verification in the desk study's figures and
\measuredcachedsavingpct{} of it in ours. Pre-computation is what makes
the desk study's own published figures fit the budget at all, which converts
a performance question into a security one: what happens when the
pre-computed pool runs out under burst?
We simulate it. The pool stalls; it does not reuse, and reuse is not
configurable, because we show precisely what reuse costs --- an observer of
two presentations sharing one blinding factor recovers the signature scalar
and every undisclosed attribute by subtraction, an extraction our test suite
performs rather than merely asserts. Stalling is not free either: under a
sustained
\poolburstratesi{} burst against a $\sim$\poolrefillratesi{} refill rate, the
worst response in a \poolwindow{} window is \poolamplification{}$\times$ an
idle one, a timing side channel with a duration-invariant signature ---
overload factor \pooloverloadfactor{}$\times$, growth per request
\poolgrowthperrequest{}$\times$ an idle response, opening at request
\poolonsetrequest{}. Separately, we ask how much of the unlinkability the
hardware layer underneath hands straight back. Partitioning a fleet
description --- a synthetic, hand-authored fixture, not telemetry from any
live fleet --- \anonfleetsize{} identical hosts, after six weeks of the
kind of operational drift in patch level and microcode such a fleet would
accumulate, fall into \anonconfigs{} distinguishable groups of
\anonsizes{}. \anonsingletons{} host is alone on its configuration before
any credential is even presented. Effective set size \anoneffectiveset{} of
\anonfleetsize{}; entropy \anonentropy{}. Two of the five layers this paper's
subject construction proposes --- hardware attestation and threshold escrow
--- are implemented here only as cost-representative stubs with no security
whatsoever, and every number that touches them says so in the sentence that
states it, not in a footnote.}

\vspace{0.6em}
\noindent{\small\textbf{Artifact.} \occultation{}~\cite{occultationtool} is
Apache-2.0. Every number in Section~\ref{sec:results} is regenerated from the
tool by \texttt{scripts/\allowbreak regen-results.sh} and read into this
document from the generated files under \texttt{../results/}; continuous
integration regenerates the host-independent half of that directory and
fails on a diff, so a stale number cannot outlive a code change. The other
half is wall-clock timing, differs by host, and is labelled as such
everywhere it appears.}

\clearpage

\tableofcontents
\clearpage

% ================================================================
\section{The Problem}
\label{sec:intro}

An enterprise sends a software agent out to buy something from a vendor. The
agent presents a credential, the vendor checks it, the order goes through.
Nobody did anything wrong.

Now ask what the vendor learned.

It learned which company the agent belongs to, because the credential said
so. It learned which employee set the agent going, because the evidence
record the standard asks for carries an \texttt{initiating\ub user} field. It
can count how many times that same agent came back this month, because every
action is bound into a hash chain whose roots are anchored where anybody can
read them. And it can line all of that up against every other agent from the
same firm, because the same identifiers run through all of them. By the end
of the quarter the vendor holds a running account of another company's
internal operations, assembled as a by-product of checking purchase orders.

The vendor had one question: is this agent allowed to do this. None of the
rest was needed to answer it.

That gap is open issue 2 in C5, the identity domain of the Proof-of-Control
standard~\cite{povocopenissues} --- six requirements, the thinnest of the
standard's six domains. Can an autonomous agent's actions be independently
verifiable and attributable when it matters, and unlinkable across relying
parties the rest of the time? A working-group member raised it in the
plainest possible terms: enterprises need to deploy agents without disclosing
customer PII or corporate identity to third-party vendors, and the best way
to secure information is not to hand it out unnecessarily.

The tension is structural, not incidental. Every other part of the standard
this domain sits inside pulls the other way, toward \emph{more} identifiable
evidence: an explicit \texttt{agent\ub id}, an \texttt{initiating\ub user}, a
hash chain that links every action to every action before it, anchors
published where anyone can see them. If you set out to design the worst
possible object for unlinkability you would land somewhere near that: a
permanent, ordered, cryptographically bound dossier. You do not fix it by
leaving a field out of a record. You need a credential layer whose
presentations do not link to each other, sitting underneath a log layer
built, by design, to chain everything together.

The cryptography for that layer is not new. The anonymous-credential line
runs back two decades, and Section~\ref{sec:related} walks it. So why is
nobody running it here? Because nobody had measured what it costs at agent
action rates. One desk study estimated it: a concrete five-layer answer ---
hardware attestation, an anonymous credential, an evidence log, a threshold
escrow for accountability, and a transport layer --- priced from published
benchmarks of other people's hardware rather than from anything it ran. We
checked the arithmetic. It contains an error large enough to flip the study's
own headline conclusion, and finding that error is most of what this paper is
about.

Here is the error. The standard allows \deskbudget{} per action. The desk
study weighed presentation against that budget: about \deskpresentnaive{},
under. Then it weighed verification: about \deskverifynaive{}, also under.
Two halves under budget, so the design fits.

An agent that presents a credential and has it verified pays for both.
\desktotalnaive{}. That is over.

No cryptographic assumption in the desk study is wrong. The mistake is in the
addition, and it survives being pointed out only because presentation and
verification happen on different machines at different times, which makes it
easy to grade them on separate report cards.

This paper reports what happens when that composed path, and the two
security questions the desk study's own numbers raise but do not answer ---
whether the fix for the composition error is itself sound, and how large an
anonymity set the hardware layer actually leaves standing --- are measured
rather than estimated.

\paragraph{One thing we are not claiming.} Two of the five layers below are
stubs with no security in them whatsoever. Our BBS+ is real cryptography that
will not interoperate with anybody else's. The fleet whose anonymity we
partition is a file we wrote by hand, not telemetry from a live one. Every
figure in this paper sits inside those limits, and each limit is repeated in
the sentence that states the number rather than parked in a footnote.

% ================================================================
\section{What Others Have Built}
\label{sec:related}

Every mechanism this design needs already exists and is trusted. Pointing
them at a machine-speed, hash-chained, publicly anchored agent log is the
part nobody has done. So every citation below does double duty: it shows the
primitive is real, and it marks the distance between that primitive and this
use of it.

\paragraph{Hardware attestation without a platform identifier.} Direct
Anonymous Attestation~\cite{daa} lets a TPM prove that \emph{some} valid TPM
produced a statement without revealing which one. It is adopted in TPM~2.0
and standardized as ISO/IEC~20008~\cite{iso20008}. The elliptic-curve
variant, ECDAA, is the practical form used in FIDO and is the closest
existing answer to the hardware layer this construction needs. It is also
the layer we do not implement: Section~\ref{sec:construction} explains why
it is modelled rather than built.

\paragraph{Anonymous credentials.} Camenisch--Lysyanskaya
credentials~\cite{cl2001}, Idemix~\cite{idemix}, and U-Prove~\cite{uprove}
are the lineage this paper's credential layer descends from. BBS+ signatures
over pairing-friendly curves~\cite{bbsplus,bls12381}, with a maturing
IETF/W3C track for selective disclosure~\cite{bbsdraft}, are the concrete
choice the desk study made and the one we implement over BLS12-381. Also
relevant, and less well known, is SPSEQ-UC~\cite{spseq}, structure-preserving
signatures on equivalence classes over updatable commitments: it targets
re-randomization directly, which is the property a log layer built from
these credentials would need and which our evidence-log layer does not yet
touch.

\paragraph{Group and ring signatures.} Chaum and van Heyst's group
signatures~\cite{chaumvanheyst} and Rivest, Shamir and Tauman's ring
signatures~\cite{ringsig} are the other route to ``some member of a group
signed this.'' Group signatures include an opening authority, which is the
shape our escrow layer generalizes to a quorum.

\paragraph{Unlinkable, bounded tokens.} Privacy Pass~\cite{privacypass} and
the VOPRF construction beneath it~\cite{voprf} are the closest deployed
analogue to ``unlinkable but rate-limited,'' and the desk study's own cost
table lists a Privacy Pass token at roughly \SI{850}{\micro\second}, sub-ms
amortized when pre-fetched in batches --- cheap enough that it is plausibly
the mechanism that makes an unlinkable design practical rather than merely
possible, though it is out of scope for the construction this paper
measures.

\paragraph{Threshold cryptography and escrow.} Threshold ElGamal, distributed
key generation, and Chaum--Pedersen DLEQ proofs for verifiable partial
decryption~\cite{chaumpedersen} are the machinery that turns ``escrow'' from
a standing backdoor into a quorum action with a public receipt. The
literature that keeps that machinery honest is the key-escrow critique in
\emph{Keys Under Doormats}~\cite{keysunderdoormats} and the specific,
documented failure modes of the Clipper Chip~\cite{clipperchip}: a
centralized depository, fragile operational control, and no transparency.
Section~\ref{sec:construction} returns to this literature when it explains
why our escrow stub's absence from the real half of this tool is not an
oversight.

\paragraph{Transport-layer linkability.} Oblivious HTTP~\cite{ohttp} and
relay-based padding schemes are easy to forget and fatal to omit: an
unlinkable credential presented over a directly attributable TCP connection
is not unlinkable. This construction, and this paper, do not touch the
transport layer at all; Section~\ref{sec:open} says so explicitly rather
than leaving it as a silent gap.

\paragraph{Where this paper sits.} None of the above measures what BBS+
costs at agent action rates, composed as a single present-and-verify path,
against a fixed real-time budget, on real hardware. That measurement --- and
the arithmetic error it corrects --- is what this paper adds. Nor does any of
it ask whether a pre-computation pool holds up under burst, which is the
security question the fix for that error raises. In most of this literature
pre-computation is a performance detail. Here it is the thing standing
between the design and its own latency budget.

% ================================================================
\section{The Construction We Measured}
\label{sec:construction}

The desk study's construction has five layers.

\begin{table}[h]
\centering\small
\caption{The five-layer construction P05 proposes. \occultation{} measures
or models rows 1, 2 and 4 (hardware attestation, agent credential,
accountability); the evidence-log and transport rows are out of scope for
this paper (Section~\ref{sec:open}).}
\label{tab:layers}
\begin{tabular}{@{}p{2.6cm}p{5.0cm}p{5.4cm}@{}}
\toprule
\textbf{Layer} & \textbf{Mechanism} & \textbf{Defeats} \\
\midrule
Hardware attestation & ECDAA: prove membership in a valid TEE group without
emitting a platform identifier & a relying party or auditor fingerprinting
the silicon \\
\addlinespace
Agent credential & BBS+ / SPSEQ-UC: fresh blinding per presentation,
selective attribute disclosure & colluding relying parties correlating
sessions \\
\addlinespace
Evidence log & ephemeral per-session chains, roots posted as
re-randomizable Pedersen commitments & a log operator reconstructing a
global action sequence \\
\addlinespace
Accountability & $k$-of-$n$ threshold ElGamal escrow tag with a NIZK that the
tag encrypts the credential-bound identity & a rogue escrow node unmasking
unilaterally \\
\addlinespace
Transport & OHTTP relay + padding & a network observer linking sessions by
IP and timing \\
\bottomrule
\end{tabular}
\end{table}

Not all five are implemented, and the three we touch are not implemented to
the same standard as each other. Table~\ref{tab:boundary} draws the line
exactly, and it is required to match \occultation{}'s \texttt{README.md}
member for member. That requirement is deliberate. A reader should be able to
check what a paper says it measured against the tool's own documentation,
rather than infer it from the fact that Table~\ref{tab:layers} labelled
something a stub --- because not every gap between a claim and the code
behind it is about a modelled component.

\begin{table}[h]
\centering\small
\caption{What \occultation{} is, stated as a three-way split, matching
\texttt{README.md} exactly. ``Real but not interoperable'' is not a hedge:
the cryptography is genuine and every property below it is asserted by a
test, but the wire encoding has not been checked against another
implementation.}
\label{tab:boundary}
\begin{tabular}{@{}p{2.8cm}p{9.9cm}@{}}
\toprule
\textbf{Category} & \textbf{Members} \\
\midrule
Real &
The Ed25519 baseline, the BLS12-381 group and pairing arithmetic, the
pre-computation pool, the anonymity-set calculator, and the timing
harness. \\
\addlinespace
Real but not interoperable &
BBS+ key generation, signing, verification, presentation and proof
verification, implemented over \texttt{blstrs} BLS12-381 against
\texttt{draft-irtf-cfrg-bbs-signatures}~\cite{bbsdraft}. Forged signatures
fail, tampered proofs fail, and undisclosed attributes stay undisclosed ---
each asserted by test. It is not validated against the draft's test
vectors, so its wire encoding will not interoperate with another BBS+
implementation. \\
\addlinespace
Modelled --- no security whatsoever &
ECDAA (\texttt{ModelledEcdaa}) and threshold ElGamal escrow
(\texttt{ModelledThresholdElGamal}). Both carry the correct interface and a
representative cost profile and provide no anonymity, no soundness, no
confidentiality and no accountability. Neither executes without an explicit
\texttt{--allow-modelled} flag; asking for one without it exits non-zero and
explains why, rather than silently running a component that provides no
security. \\
\bottomrule
\end{tabular}
\end{table}

Three of the five rows are what the rest of this paper touches: hardware
attestation, the agent credential, and accountability. It is worth being
blunt about how little of Section~\ref{sec:results} is a timing on this
machine. Exactly one table is: Table~\ref{tab:primitives}, the
agent-credential row, real BLS12-381 pairing arithmetic and a real BBS+
implementation on real hardware. Nothing else in the section is.
Table~\ref{tab:composed}'s desk-study columns are the published study's own
constants, never run here. Table~\ref{tab:anonymity} partitions a fixed TOML
file, which has no cryptography in it to time.
Section~\ref{sec:results-pool}'s burst figures come out of a simulation fed
operator-supplied costs. And every figure that touches ECDAA or threshold
escrow --- the hardware-attestation and accountability rows --- is a cost
estimate produced by a stub with no cryptographic content in it. We use
neither stub to support any claim about security. Only, where unavoidable,
about cost, and always labelled.

Leaving escrow out of the real half of Table~\ref{tab:boundary} is a choice,
not an oversight, and Section~\ref{sec:related}'s escrow literature is the
reason. A $k$-of-$n$ threshold ElGamal tag with a NIZK of correct encryption
is exactly the machinery \emph{Keys Under Doormats}~\cite{keysunderdoormats}
and the Clipper Chip's documented failure~\cite{clipperchip} warn is easy to
build and hard to build safely: the danger sits in the governance around a
working implementation as much as in the cryptography. A stub built only to
carry a representative cost cannot be mistaken for a governance guarantee we
never evaluated. Building a real threshold-opening path and asserting its
security is future work. It is not something a cost-measurement tool should
do as a side effect.

The evidence-log layer of Table~\ref{tab:layers} is not implemented at all.
It is P05's open questions 1 and 2 (Section~\ref{sec:open}), and nothing we
measured bears on it. The transport layer is not implemented either, but it
is a separate omission rather than one of those two questions, and
Section~\ref{sec:open} names it on its own instead of folding it in.

% ================================================================
\section{Measuring the Composed Path}
\label{sec:method}

The desk study's mistake was structural. It had one budget and two numbers,
and it held each number up against the budget on its own.
\occultation{}'s \texttt{bench --composed} exists to make that move
unavailable. It folds a presentation cost and a verification cost into one
\texttt{Composed} value, adds them, and judges the sum --- never a half ---
against the \deskbudget{} budget. That budget is written down once, in
\texttt{src/cost.rs}, and \texttt{bench} is the only command that reads it;
\texttt{pool} and \texttt{anonymity} answer different questions and never
consult it. A verdict has three values rather than two: \texttt{Fits},
\texttt{Misses}, and the one that earns its place, \texttt{Marginal}, for a
composition under budget by no more than 5\,\% of it. Drop that third value
and \desktotalcached{} against a \deskbudget{} budget prints as a pass ---
which is the desk study's own verdict on its cached-pairing figure, and the
verdict this tool exists to refuse.

\texttt{bench --composed} prints two profiles side by side and never merges
them. The \emph{desk-study} profile repeats the published figures verbatim,
tagged \texttt{PUBLISHED}, with the citation attached: constants, not
measured on this stack, reported so that anybody can check the composition
arithmetic without borrowing our hardware. The \emph{measured} profile times
the same operations here, seeded RNG, median of many iterations, tagged
\texttt{measured here}. They are supposed to disagree in magnitude --- they
came off different silicon --- and Section~\ref{sec:results} says by how
much and in what shape.

One row of the measured profile carries a third tag, \texttt{ESTIMATED}, and
what it is not needs saying precisely. Nothing in \occultation{} measures a
pooled presentation's \emph{online} phase --- the work left over once a
blinding factor has been pre-computed and pulled off the pool. That row is
derived from a cheaper stand-in: two hash-to-scalar operations, standing in
for the challenge hash and the field-arithmetic responses a real online phase
would do. \texttt{occultation pool}, the command that actually drives the
pool under load (Section~\ref{sec:pool}), does not close the gap either ---
when the operator supplies no \texttt{--online} cost, it defaults to this
same estimate. So the estimate and the pool simulation are not independent
evidence for each other. And the row's own note says the estimate is almost
certainly too low: it implies a \pooledreductionpct{} reduction in
presentation cost where the desk study's own pooled figure implies
\deskpooledreductionpct{}. We print it anyway, because dropping the row would
leave the measured profile's pooled composition unstated altogether, and a
caveated estimate beats a silent gap. It is not a claim that we measured a
pooled presentation. No number in this paper should be read as one.

% ================================================================
\section{What It Actually Costs}
\label{sec:results}

Every table and every inline figure in this section is generated by
\texttt{scripts/\allowbreak regen-results.sh} from files under
\texttt{../results/} and either \texttt{\textbackslash{}input}-ed or
substituted as a \LaTeX{} macro. None of it was typed. Files under
\texttt{results/deterministic/} are arithmetic over the desk study's
published constants, a seeded pool simulation (\texttt{pool-facts.tex}
included), or a partition of a fixed fleet file, so they come out
byte-for-byte identical on any host and continuous integration fails the
build on a diff. Files under \texttt{results/measured/} are wall-clock
timing. They differ by host, so they are not diffed, and every number taken
from them is marked as measured --- in its own table, or, for
\texttt{composed-facts.tex}, in the prose that quotes it.

\subsection{Adding the two halves together}
\label{sec:results-composed}

% Generated by scripts/regen-results.sh from the desk_study rows of
% results/measured/composed.json. Those rows are the desk study's published
% figures, which are constants and do not vary by host, so this file lives
% under deterministic/ and is diffed by CI. Do not hand-edit.
\begin{table}[h]
\centering\small
\begin{tabular}{>{\raggedright\arraybackslash}p{4.3cm}rrr>{\raggedright\arraybackslash}p{3.7cm}}
\toprule
Composition & Present & Verify & Total & Verdict \\
\midrule
naive (no pool, no cached pairings) & 13.80\,ms & 4.20\,ms & 18.00\,ms & MISSES by 3.00\,ms \\
cached pairings only & 13.80\,ms & 1.10\,ms & 14.90\,ms & MARGINAL --- 0.10\,ms spare \\
pre-computation pool + cached pairings & 2.80\,ms & 1.10\,ms & 3.90\,ms & fits, 11.10\,ms spare \\
\bottomrule
\end{tabular}
\caption{The composed present-and-verify path against the 15.00\,ms budget: the desk study's published figures for presentation and verification, assessed together rather than separately as the desk study assessed them. The naive composition sums to 18.00\,ms and misses the budget by 3.00\,ms even though each half alone appeared to fit. Caching pairings brings the sum to 14.90\,ms, 0.10\,ms under the limit --- not headroom, hence \texttt{MARGINAL} rather than a pass. Only the pre-computation pool, combined with cached pairings, fits with room to spare: 3.90\,ms total, 11.10\,ms of headroom. \texttt{results/\allowbreak measured/\allowbreak composed.txt} reports the profile measured on this host, which diverges from these published figures by roughly an order of magnitude without changing the correction, which is an error in arithmetic, not in hardware.}
\label{tab:composed}
\end{table}


Table~\ref{tab:composed} is the desk study's own published figures, composed
as one path instead of two. The naive composition misses the budget outright.
Caching pairings pulls it back to \texttt{MARGINAL} rather than a pass ---
under the limit, but by less than the jitter of an uncontended context
switch. Only the pre-computation pool, together with cached pairings, fits
with room. That is the headline correction, and it is a statement about these
figures: in the desk study's own numbers, pre-computation is not an
optimization sitting on top of a design that already works. It is the reason
the design fits its budget at all.

Table~\ref{tab:primitives} times the individual primitives underneath
Table~\ref{tab:composed}'s presentation and verification costs, here on this
host rather than read out of a publication: ten primitive operations instead
of three compositions, and no budget verdict against any one of them, which
is the whole point Table~\ref{tab:composed} exists to make.

% Generated by scripts/regen-results.sh from results/measured/primitives.json.
% Wall-clock timing: differs per host, not diffed by CI. See
% results/measured/HEADER.txt for the host and toolchain that produced this
% file. Do not hand-edit.
\begin{table}[h]
\centering\small
\begin{tabular}{lrrrl}
\toprule
Operation & Cost & P95 & $\times$Ed25519 & Provenance \\
\midrule
Ed25519 sign & 0.007\,ms & 0.007\,ms & 0.3 & measured here (n=200) \\
Ed25519 verify & 0.018\,ms & 0.019\,ms & 0.7 & measured here (n=200) \\
Ed25519 sign + verify (baseline) & 0.025\,ms & -- & 1.0 & measured here (n=200) \\
G1 scalar multiplication & 0.039\,ms & 0.040\,ms & 1.6 & measured here (n=200) \\
hash\_to\_scalar & 0.002\,ms & 0.002\,ms & 0.1 & measured here (n=200) \\
BBS+ sign (issuer) & 0.208\,ms & 0.219\,ms & 8.2 & measured here (n=200) \\
BBS+ signature verify (holder) & 0.642\,ms & 0.698\,ms & 25.3 & measured here (n=200) \\
BBS+ presentation & 0.539\,ms & 0.564\,ms & 21.2 & measured here (n=200) \\
BBS+ proof verify & 0.766\,ms & 0.793\,ms & 30.2 & measured here (n=200) \\
BBS+ proof verify, cached pairings & 0.684\,ms & 0.713\,ms & 27.0 & measured here (n=200) \\
\bottomrule
\end{tabular}
\caption{Primitive costs measured on this host against an Ed25519 sign-and-verify round trip. ECDAA and threshold escrow are omitted: they are modelled stubs and this table is generated from a plain \texttt{occultation bench} run, which never executes a modelled component. See \texttt{results/measured/HEADER.txt} for the exact host and toolchain.}
\label{tab:primitives}
\end{table}


The measured composed path (\texttt{results/measured/composed.txt}, given in
prose here rather than as a table) is \measureddivergencefactor{}$\times$
\measureddivergencedirection{} here than the desk study's published figures
imply, and the composed BBS+ path costs \measuredratiotobaseline{}$\times$ an
Ed25519 sign-and-verify round trip on this host. That factor is not
like-for-like. The desk study's naive presentation figure folds in a
\SI{4.6}{\milli\second} (33\,\%) allowance for the escrow tag; the measured
figure here leaves escrow out entirely. Take escrow off both sides and the
divergence is roughly \measureddivergencefactorexclescrow{}$\times$. None of
which touches the composition error, because that error is arithmetic rather
than a property of any one machine, and it comes out the same wherever the
desk study's published figures are added up.

The magnitude is the part that hides the finding. The divergence is not
uniform. Cached pairings are the one optimization measured on both sides:
the desk study's figures credit caching the pairing computation with
\deskcachedsavingpct{} of verification cost, and the same optimization saves
\measuredcachedsavingpct{} here. A host that were simply faster across the
board would keep the desk study's ratio between optimized and unoptimized
verification. This host does not. So the measured block does not corroborate
the desk study's optimization structure, and we report that as a finding in
its own right rather than averaging it away as noise.

\subsection{What the pool costs when it runs dry}
\label{sec:results-pool}

Section~\ref{sec:pool} argues that what the correction above really raises is
the pool's soundness, not its speed. These are the numbers behind that
argument. They come from the shipped burst configuration --- capacity
\poolcapacity{}, refill \poolrefillcostsi{}, online \poolonlinecostsi{},
burst rate \poolburstratesi{}, seed \poolseed{}, the same configuration
\texttt{tests/acceptance.rs}'s acceptance test~2 exercises. Of
\poolrequests{} requests over a \poolwindow{} window, \poolserved{} were
served immediately and \poolstalled{} stalled. The worst response was
\poolamplification{}$\times$ an idle one.

That last figure is a peak over the observation window, and it grows linearly
with the window, so it should never travel on its own. Beside it go the
three numbers that do not depend on how long you watch: an overload factor of
\pooloverloadfactor{}$\times$, arrivals over refill; a growth rate of
\poolgrowthperrequest{}$\times$ an idle response, added to each successive
stall once the pool is dry; and an onset at request \poolonsetrequest{},
which is how briefly a burst has to be sustained before the channel opens.
Every one of \pooldistinctblinding{} requests got a distinct blinding factor.
Reuse events: \poolreuseevents{}.

The refill and online costs above (\poolrefillcostsi{} and
\poolonlinecostsi{}) are values \texttt{scripts/regen-results.sh} stipulates
for this simulation. They are not a measurement of any deployment. Both
\poolamplification{}$\times$ and \poolgrowthperrequest{}$\times$ scale
inversely with whatever online cost you pick; \pooloverloadfactor{}$\times$
does not depend on it at all.

\subsection{How much of the fleet looks alike}
\label{sec:results-anonymity}

% Generated by scripts/regen-results.sh from
% results/deterministic/anonymity-drifted.json, a partition of the fixed
% fleet file examples/fleet-drifted.toml. Do not hand-edit.
\begin{table}[h]
\centering\small
\begin{tabular}{rlccc}
\toprule
Size & Host group & PCE SVN & Microcode & TDX module \\
\midrule
71 & us-east-pool-patched & 13 & 0x2b000603 & 1.5.05 \\
34 & us-east-pool-pending & 12 & 0x2b000603 & 1.5.05 \\
11 & eu-west-rack-7 & 12 & 0x2b000401 & 1.5.03 \\
3 & us-east-canary & 13 & 0x2b000701 & 1.5.05 \\
1 & us-east-rebuilt-0419 & 14 & 0x2b000701 & 1.5.06 \\
\bottomrule
\end{tabular}
\caption{Anonymity sets for fleet \texttt{acme-prod-drifted}: 120 hosts of identical hardware, after six weeks of ordinary operational drift in patch level, microcode and PCE SVN. 5 distinct TCB configurations emerge; the largest anonymity set is 71 of 120 hosts, and 1 host(s) are alone on their configuration --- no credential scheme can help them, because the quote collateral identifies the machine before any credential is even presented. Effective set size 52.73 of 120; entropy 1.470 bits.}
\label{tab:anonymity}
\end{table}


Table~\ref{tab:anonymity} answers a question the desk study's cost analysis
never asks. Suppose the credential layer is perfectly unlinkable. How large
is the anonymity set the hardware attestation layer leaves standing
underneath it? \texttt{examples/fleet-drifted.toml} is a synthetic,
hand-authored fixture, not telemetry from any live fleet. It describes
\anonfleetsize{} hosts of identical hardware after six weeks of the kind of
operational drift a real fleet accumulates: a canary batch here, a delayed
patch there, a microcode revision that reached some racks and not others. A
uniform fleet with none of that drift
(\texttt{examples/fleet-uniform.toml}) is one anonymity set the size of the
whole fleet, by construction, and exists only as the tool's control case.
Table~\ref{tab:anonymity} is the drifted fleet, which is closer to how a
fleet actually runs --- though both files are things we wrote, not telemetry
from a live fleet.

% ================================================================
\section{Is the Pre-Computation Pool Sound?}
\label{sec:pool}

Section~\ref{sec:results-composed} established that in the desk study's own
published figures, pre-computation is what lets the composed path fit its
budget at all, not an optional speed-up on top of a design that already
worked. That moves the risk somewhere else. The question stops being whether
BBS+ is fast enough and becomes whether the pool that makes it fast enough is
sound. Nobody asked before, because until the arithmetic error surfaced,
pre-computation looked like an optimization, and optimizations are not where
a security review looks.

\paragraph{The two available behaviours.} A pool of pre-computed blinding
factors, sitting in front of a burst of demand, has exactly two things it can
do once it is empty. It can make the caller wait for a fresh item, or it can
hand back one that has already been used. \occultation{}'s
\texttt{PrecomputationPool} does the first and never the second. On
exhaustion \texttt{acquire} returns \texttt{None} rather than a used item,
and there is no flag and no code path that makes it return one.
Section~\ref{sec:results-pool} is the price of that choice: under a sustained
burst the delay does not plateau, it grows with the length of the burst, and
that growth is a timing side channel with a duration-invariant signature.
Watch nothing but an agent's response latency and you can read off how
overloaded it is and how long it has been that way.

\paragraph{Why the alternative is not a lesser evil.} Reusing one set of
blinding scalars across two presentations does not merely weaken
unlinkability. It discloses the credential outright, and that is provable
rather than asserted. A BBS+ presentation's proof of knowledge has the
Schnorr-style shape common to sigma protocols: for a hidden value $x$ (the
signature's blinding scalar, or an undisclosed attribute) with blinding
$\tilde{x}$ and Fiat--Shamir challenge $c$, the response the verifier sees
is
\begin{equation}
\hat{x} = \tilde{x} + c \cdot x \pmod{r}.
\end{equation}
If two presentations to two different relying parties reuse the same
blinding $\tilde{x}$ but necessarily bind to two different challenges $c_1
\neq c_2$ --- the challenge is derived from a per-session nonce, so
identical challenges would themselves require a coincidence --- then
subtracting the two responses cancels the blinding entirely:
\begin{equation}
\hat{x}_1 - \hat{x}_2 = (c_1 - c_2)\, x
\quad\Longrightarrow\quad
x = \frac{\hat{x}_1 - \hat{x}_2}{c_1 - c_2}.
\end{equation}
This is not a hypothetical weakness. \texttt{pool.rs}'s test
\texttt{reusing\ub one\ub set\ub of\ub scalars\ub leaks\ub the\ub signature\ub and\ub every\ub hidden\ub attribute}
does exactly this division on the reused presentation's responses, against a
real signature and real hidden attributes, and asserts that what comes back
is the signature's scalar and every undisclosed attribute in order. The two
presentations are linkable on sight before any of that arithmetic is done
anyway, because they carry the same commitment values. Reuse is not a milder
version of unlinkability failing. It is unlinkability's negation, in closed
form, handed to anyone who sees both transcripts.

\paragraph{What the type system buys, stated precisely.} \texttt{PoolItem}
is not \texttt{Clone}, and its \texttt{consume} method takes \texttt{self} by
value, so \emph{accidental} reuse --- holding an item, copying it, or pulling
it off the pool a second time --- is impossible. The fields are private, the
type cannot be duplicated, and \texttt{acquire} never puts an item back in
its own queue. This is worth stating exactly, because it is easy to
overstate. \emph{Deliberate} reuse is not prevented. \texttt{consume} puts no
bound on what its closure returns, the scalars underneath are an ordinary
\texttt{Clone}able type with public fields, so \texttt{item.consume(|s|
s.clone())} yields an owned copy that outlives the item --- and the
lower-level proving function that takes pre-chosen scalars is public too.
What the design buys is not a guarantee that reuse cannot happen. It is that
any reuse is one conspicuous, greppable \texttt{.clone()} at the point it
happens, rather than something a pool running dry under load produces on its
own --- which is exactly the failure mode this paper's measurement is about.
A reviewer looking for that call will find it. A load test will not, because
nothing under load produces it. Exhaustion produces a stall, and only a
stall.

\paragraph{What we are not claiming.} We have not built or evaluated any
third option beyond stall and reuse. P05's own open question --- whether pool
refill, running as idle-time work inside the protected environment,
correlates with agent activity closely enough that an observable host
resource (power draw, cache behaviour, scheduling) gives it away --- is
untouched by anything we measured. And \texttt{pool}'s simulation treats
refill as a fluid approximation of a single worker rather than a
discrete-event queue. The approximation is written down in the code and it
does not change which of the two behaviours gets chosen, but it does mean the
specific latency numbers in Section~\ref{sec:results-pool} are a model of one
refill strategy rather than a measurement of every plausible one.

% ================================================================
\section{What We Still Don't Know}
\label{sec:open}

Four questions this tool does not answer, in the order the desk study's own
open-questions list puts them:

\begin{enumerate}[itemsep=3pt,topsep=3pt]

\item \textbf{Can a log operator correlate re-randomized commitments by
posting time and size, even if the commitments themselves hide everything
cryptographically?} The construction's evidence-log layer proposes
per-session ephemeral chains with roots posted as re-randomizable Pedersen
commitments, specifically so that a log operator cannot reconstruct a global
action sequence out of what the commitments contain. Nothing in this tool
touches that layer --- it is not implemented, measured, or modelled here.
Hiding a value cryptographically says nothing about when and how often the
value gets posted, and an operator who can watch posting times and sizes has
a side channel no property of the commitment scheme addresses. This is the
gap P05 names, and the thing to prove or refute. This paper does neither.

\item \textbf{What does tamper-evidence mean once the global chain is
replaced by independent per-session chains?} A single global hash chain
proves that no action was deleted. Independent per-session chains, the
construction's own proposal for unlinkability, prove only that no action was
deleted \emph{from a session an auditor already knows exists}: an agent that
suppresses an entire session leaves no gap for anyone to notice. This may be
the more consequential result available from this whole research direction
--- that unlinkability and completeness-of-record are in tension rather than
independently achievable --- and this paper's measurements say nothing about
it either way, because the evidence-log layer is out of scope here.

\item \textbf{Revocation.} Revoking a compromised agent's credential without
linking its past actions, and without an accumulator update cost that scales
badly with fleet size, is historically the hardest part of deploying
anonymous credentials, and it is barely touched by the desk study this paper
corrects. \occultation{} has no revocation mechanism, measured or modelled.

\item \textbf{Is an anonymity set built out of TCB configuration diversity
even the right set to be measuring?} Section~\ref{sec:results-anonymity}
answers ``how large is the set a fleet's hardware diversity leaves
standing,'' which is worth answering on its own. It is not the same question
as ``what anonymity set should a colluding pair of relying parties be
evaluated against.'' The right set might be agents inside one enterprise, or
enterprises sharing one credential issuer, or TEE hardware generations across
many unrelated fleets --- and the answer changes who the scheme protects from
whom. Our calculator partitions whatever fleet description you hand it. It
has no opinion about which fleet is the relevant one. Supplying that opinion
is future work, not something this measurement settles on its own.

\end{enumerate}

Two structural omissions sit above all four and are worth naming out loud.
The evidence-log and transport layers of Table~\ref{tab:layers} are not
implemented, measured, or modelled anywhere in this artifact. The first is
questions 1 and 2 above. The second --- Oblivious HTTP, relay padding, and
everything else that stops an unlinkable credential from riding over a
directly attributable connection --- is named by P05 and untouched by this
tool entirely. Every figure this paper reports, about cost or exhaustion or
anonymity, is a statement about the credential and hardware-attestation
layers alone.

% ================================================================
\section{Conclusion}

Two numbers were held up against one budget, one at a time, and each of them
cleared it. The design they belong to was reported as fitting. It does not
fit: the sum is over. That is the correction, it needs no hardware to
reproduce, and it comes out the same wherever the desk study's published
figures are added up. What it changes is what the design is standing on. In
those figures the pre-computation pool is not a speed-up bolted onto
something that already worked --- it is the reason anything fits, which turns
a performance detail into a security question nobody had reason to ask. So we
asked it. The pool stalls rather than reusing, because reuse hands anybody
who sees two transcripts the signature scalar and every undisclosed attribute
in closed form; and stalling is not free either, because the delay grows with
the burst and an observer can read the overload off an agent's response
latency.

The part of this we are least able to defend is the evidence underneath it.
The fleet whose anonymity sets we partition is a TOML file we wrote by hand,
and the drift in it is the drift we thought a real fleet would accumulate,
not drift anybody observed. Two of the five layers are stubs with no security
in them whatsoever. The BBS+ under the credential layer is real cryptography,
but it has never been run against the IETF draft's test vectors and will not
interoperate with another implementation. An error in somebody else's
arithmetic is cheap to be right about, because the arithmetic is public and
anyone can redo it. Everything else here rests on fixtures we wrote for
ourselves. If one sentence survives this paper, we would rather it be that
one than the correction.

% ================================================================
\bibliographystyle{unsrt}
\bibliography{references}

\end{document}
